\documentclass[11pt]{article}
\usepackage[utf8]{inputenc}
\usepackage{amsmath,amssymb}
\usepackage[margin=1in]{geometry}
\usepackage{hyperref}
\emergencystretch=3em

\title{Perturbations divisible by $u^k$ do not change the leading term:\\
$Z(a+u^k\tilde J)\sim Z(a)\sim C\,n^{-p/2}(\log n)^m$}
\author{Claude, with Daniel Murfet}
\date{2026-09-06}

\begin{document}
\maketitle
\sloppy

\begin{abstract}
The first multiplicity result for non-constant $\xi$. If $\xi(u)=a+u^k\tilde J(u)$ with $\tilde J$
continuous on the box and $|\tilde J|\le L$, then with $t=\sqrt n\,u^k$ the exponent is
$\beta ta+\beta t^2\tilde J/\sqrt n$ and the one-term exponential remainder bound gives, for
$n\ge16L^2$, $|Z(\xi)-Z(a)|\le\frac{4L}{\sqrt n}Z_{\beta/2,2a}(n)$. The right-hand side has the same
order $n^{-p/2}(\log n)^m$ as $Z(a)$ by unit~59 at $(\beta/2,2a)$, so $Z(\xi)\sim Z(a)$ and the
multiplicity theorem holds verbatim for such $\xi$. Zero \texttt{sorry}/\texttt{axiom}.
\end{abstract}

\section{Setting}
The Taylor tree's higher layers are not individually lower order in general (a perturbation $J$ that
does not vanish on a minimal divisor component contributes at leading order). Divisibility by $u^k$ is
the natural sufficient condition: every layer then carries an extra $n^{-p/2}$, and the whole tail is
controlled at once by $|e^x-1|\le|x|e^{|x|}$ together with $xe^{-x}\le1$.

\section{The things formalised}
{\small
\begin{verbatim}
theorem exp_pert_abs_le (β a L ε t y : ℝ) (hβ : 0 < β) (hL : 0 ≤ L) (hε : 0 ≤ ε)
    (hy : |y| ≤ L) (hLε : L * ε ≤ 1 / 4) :
    |Real.exp (-β * t ^ 2 + β * t * (a + t * y * ε)) - Real.exp (-β * t ^ 2 + β * t * a)|
      ≤ 4 * L * ε * Real.exp (-(β / 2) * t ^ 2 + β * t * a)
theorem boxIntegralXi_div_sub_le ... (hnL : 16 * L ^ 2 ≤ n) :
    |boxIntegralXi β a b (Real.sqrt n) k h (fun u => (∏ i, u i ^ k i) * J u)
        - boxIntegralFin β a b (Real.sqrt n) k h|
      ≤ 4 * L / Real.sqrt n * boxIntegralFin (β / 2) (2 * a) b (Real.sqrt n) k h
theorem boxIntegralXi_div_isEquivalent ...
    (J : (Fin (D + 2) → ℝ) → ℝ) (hJ : Continuous J) :
    ∃ (p : ℝ) (m : ℕ) (C : ℝ), 0 < C ∧ (∀ i, p ≤ finExp k h i) ∧
      m + 1 = (Finset.univ.filter fun i => finExp k h i = p).card ∧
      (fun n => boxIntegralXi β a b (Real.sqrt n) k h (fun u => (∏ i, u i ^ k i) * J u))
        ~[atTop] fun n => C * (n ^ (-(p / 2)) * Real.log n ^ m)
\end{verbatim}
}

\section{Proof strategy}
Pointwise, $e^{-\beta t^2+\beta ta}|e^{\beta t^2y\varepsilon}-1|\le e^{-\beta t^2+\beta ta}\,
\beta t^2L\varepsilon\,e^{\beta t^2/4}=4L\varepsilon\,e^{-(\beta/2)t^2+\beta ta}\cdot
\tfrac{\beta t^2}{4}e^{-\beta t^2/4}$ and the last factor is at most $1$
(\texttt{Real.add\_one\_le\_exp}). Integrating over the box (integrands continuous, hence
integrable; the pointwise bound holds almost everywhere by \texttt{ae\_restrict\_mem}) gives the
perturbation bound. For the asymptotic, unit~59 is applied at $(\beta,a)$ and at $(\beta/2,2a)$; the
two minima and multiplicities coincide (both are read off from $(k,h)$), the difference is
$O(n^{-1/2}Z_{\beta/2,2a})=o(Z(a))$ via \texttt{IsLittleO.mul\_isBigO} with
$n^{-1/2}=o(1)$, and \texttt{IsEquivalent.add\_isLittleO} concludes.

\section{Roadblocks and resolutions}
\begin{itemize}
\item A \texttt{set} of $t=\sqrt n\prod u^k$ before unfolding the integrands did not fold the
later-introduced occurrences, and a subsequent rewrite of $\prod u^k$ hit the inside of $t$; introduce
$t$ as an opaque variable after unfolding and rewrite $\sqrt n\prod u^k$ into it.
\item Nonnegativity of the box integral needs the almost-everywhere box membership
(\texttt{integral\_nonneg\_of\_ae}), since $\prod u_i^{h_i}$ is not globally nonnegative.
\item \texttt{isLittleO\_one\_iff} takes the codomain type explicitly.
\end{itemize}

\section{What was Mathlib, what was new}
Mathlib: \texttt{Real.add\_one\_le\_exp}, \texttt{Real.exp\_le\_exp}, \texttt{norm\_integral\_le\_of\_norm\_le},
\texttt{integral\_nonneg\_of\_ae}, \texttt{isLittleO\_one\_iff}, \texttt{tendsto\_rpow\_neg\_atTop},
\texttt{IsLittleO.mul\_isBigO}, \texttt{IsEquivalent.add\_isLittleO}. New: the perturbation bound and
the multiplicity theorem for $\xi=a+u^k\tilde J$.

\section{Lessons}
Comparing $Z(\xi)$ with a constant-$\xi$ integral at a \emph{different} temperature $(\beta/2,2a)$ is
the trick that makes the remainder uniform: the extra $e^{\beta t^2/4}$ from the perturbation is
absorbed into the kernel, and the asymptotic theory for constant $\xi$ is available at every
temperature.

\section{Outcome}
The multiplicity theorem now covers a genuine class of non-constant $\xi$: those whose variation is
divisible by the resolution monomial $u^k$. Perturbations that do not vanish on the minimal divisor
components remain part of the leading coefficient (the state-density picture of the paper).
\end{document}
